\begin{lemma}
\label{lemma-closed-point-nearby-fibre}
Let $f : X \to S$ be a morphism of schemes.
Let $x \leadsto x'$ be a specialization of points in $X$.
Set $s = f(x)$ and $s' = f(x')$.
Assume
\begin{enumerate}
\item $x'$ is a closed point of $X_{s'}$, and
\item $f$ is locally of finite type.
\end{enumerate}
Then the set
$$
\{x_1 \in X
\text{ such that }
f(x_1) = s
\text{ and }
x_1\text{ is closed in }X_s
\text{ and }
x \leadsto x_1 \leadsto x'
\}
$$
is dense in the closure of $x$ in $X_s$.
\end{lemma}

\begin{proof}
We apply
Schemes, Lemma \ref{schemes-lemma-points-specialize}
to the specialization $x \leadsto x'$.
This produces a morphism $\varphi : \Spec(B) \to X$
where $B$ is a valuation ring such that $\varphi$ maps the
generic point to $x$ and the closed point to $x'$. We may also
assume that $\kappa(x)$ is the fraction field of $B$.
Let $A = B \cap \kappa(s)$. Note that this is a valuation ring (see
Algebra, Lemma \ref{algebra-lemma-valuation-ring-cap-field})
which dominates the image of $\mathcal{O}_{S, s'} \to \kappa(s)$.
Consider the commutative diagram
$$
\xymatrix{
\Spec(B) \ar[rd] \ar[r] &
X_A \ar[d] \ar[r] & X \ar[d] \\
& \Spec(A) \ar[r] & S
}
$$
The generic (resp.\ closed) point of $B$ maps to a point $x_A$
(resp.\ $x'_A$) of $X_A$ lying over the generic (resp.\ closed)
point of $\Spec(A)$. Note that $x'_A$ is a closed point
of the special fibre of $X_A$ by
Morphisms,
Lemma \ref{morphisms-lemma-base-change-closed-point-fibre-locally-finite-type}.
Note that the generic fibre of $X_A \to \Spec(A)$ is isomorphic
to $X_s$. Thus we have reduced the lemma to the case where $S$ is
the spectrum of a valuation ring, $s = \eta \in S$ is the generic point, and
$s' \in S$ is the closed point.

\medskip\noindent
We will prove the lemma by induction on $\dim_x(X_\eta)$.
If $\dim_x(X_\eta) = 0$, then there are no other points of $X_\eta$
specializing to $x$ and $x$ is closed in its fibre, see
Morphisms, Lemma \ref{morphisms-lemma-quasi-finite-at-point-characterize},
and the result holds. Assume $\dim_x(X_\eta) > 0$.

\medskip\noindent
Let $X' \subset X$ be the reduced induced scheme structure on
the irreducible closed subscheme $\overline{\{x\}}$ of $X$, see
Schemes, Definition \ref{schemes-definition-reduced-induced-scheme}.
To prove the lemma we may replace $X$ by $X'$ as this only decreases
$\dim_x(X_\eta)$. Hence we may also assume that $X$ is an integral scheme
and that $x$ is its generic point. In addition, we may replace $X$ by an
affine neighbourhood of $x'$. Thus we have $X = \Spec(B)$ where
$A \subset B$ is a finite type extension of domains. Note that in
this case $\dim_x(X_\eta) = \dim(X_\eta) = \dim(X_{s'})$, and that in fact
$X_{s'}$ is equidimensional, see
Algebra,
Lemma \ref{algebra-lemma-finite-type-domain-over-valuation-ring-dim-fibres}.

\medskip\noindent
Let $W \subset X_\eta$ be a proper closed subset (this is the
subset we want to ``avoid''). As $X_s$ is of finite type over a field
we see that $W$ has finitely many irreducible components
$W = W_1 \cup \ldots \cup W_n$. Let
$\mathfrak q_j \subset B$, $j = 1, \ldots, r$
be the corresponding prime ideals. Let $\mathfrak q \subset B$
be the maximal ideal corresponding to the point $x'$.
Let $\mathfrak p_1, \ldots, \mathfrak p_s \subset B$ be the
minimal primes lying over $\mathfrak m_AB$. There are finitely
many as these correspond to the irreducible components of the
Noetherian scheme $X_{s'}$. Moreover, each of these irreducible
components has dimension $> 0$ (see above) hence we see that
$\mathfrak p_i \not = \mathfrak q$ for all $i$.
Now, pick an element $g \in \mathfrak q$ such that
$g \not \in \mathfrak q_j$ for all $j$ and $g \not \in \mathfrak p_i$
for all $i$, see
Algebra, Lemma \ref{algebra-lemma-silly}.
Denote $Z \subset X$ the locally principal closed subscheme defined by $g$.
Let $Z_\eta = Z_{1, \eta} \cup \ldots \cup Z_{n, \eta}$, $n \geq 0$
be the decomposition of the generic fibre of $Z$ into irreducible
components (finitely many as the generic fibre is Noetherian).
Denote $Z_i \subset X$ the closure of $Z_{i, \eta}$.
After replacing $X$ by a smaller affine neighbourhood
we may assume that $x' \in Z_i$ for each $i = 1, \ldots, n$.
By construction $Z \cap X_{s'}$ does not contain any irreducible
component of $X_{s'}$. Hence by
Lemma \ref{lemma-locally-principal-vertical}
we conclude that $Z_\eta \not = \emptyset$! In other words
$n \geq 1$. Letting $x_1 \in Z_1$ be the generic point we see
that $x_1 \leadsto x'$ and $f(x_1) = \eta$.
Also, by construction $Z_{1, \eta} \cap W_j \subset W_j$
is a proper closed subset. Hence every irreducible component of
$Z_{1, \eta} \cap W_j$ has codimension $\geq 2$ in $X_\eta$
whereas $\text{codim}(Z_{1, \eta}, X_\eta) = 1$ by
Algebra, Lemma \ref{algebra-lemma-minimal-over-1}.
Thus $W \cap Z_{1, \eta}$ is a proper closed subset.
At this point we see that the induction hypothesis applies to
$Z_1 \to S$ and the specialization $x_1 \leadsto x'$.
This produces a closed point $x_2$ of $Z_{1, \eta}$ not contained
in $W$ which specializes to $x'$. Thus we obtain
$x \leadsto x_2 \leadsto x'$, the point $x_2$ is closed in $X_\eta$,
and $x_2 \not \in W$ as desired.
\end{proof}